\section{Results}

\begin{table*}[t]
\centering
\caption{Illustrative aggregate results. Mean $\pm$ standard error over simulation seeds;
hardware success is reported over 48 trials. The best value is bold and the second best is
underlined. All planners use the same locomotion controller and sensing stream.}
\label{tab:main}
\small
\setlength{\tabcolsep}{7pt}
\begin{tabular}{lcccccc}
\toprule
Method & Sim. success $\uparrow$ & HW success $\uparrow$ & Interventions $\downarrow$
& Speed ratio $\uparrow$ & Height MAE (cm) $\downarrow$ & ECE $\downarrow$ \\
\midrule
Reactive     & $71.4\!\pm\!1.8$ & 31/48 & 17.1\% & \best{0.97} & 7.8 & --- \\
DecayMap     & $79.6\!\pm\!1.5$ & 35/48 & 13.5\% & 0.90 & 6.1 & 0.142 \\
Memory-only  & $84.8\!\pm\!1.3$ & 39/48 & 10.9\% & \second{0.94} & \second{4.7} & --- \\
\method{}    & \best{$92.1\!\pm\!0.9$} & \best{44/48} & \best{7.5\%}
             & 0.92 & \best{4.1} & \best{0.036} \\
\bottomrule
\end{tabular}
\end{table*}

\subsection{Navigation Reliability}

Table~\ref{tab:main} shows that geometry memory alone improves simulated success by 13.4
points over Reactive, confirming the value of evidence across occlusions. Adding calibrated
risk produces a further 7.3-point gain while reducing normalized speed by only 0.02 relative
to Memory-only. On hardware, \method{} completes 44 of 48 trials. Its four failures consist
of two correlated map-warp errors, one unmodeled compliant obstacle, and one state-estimator
reset. Figure~\ref{fig:interventions} shows the largest benefit under structured dropout,
where old terrain can appear plausible for several steps.

A paired bootstrap over courses gives a $-6.1$ percentage-point intervention difference
between \method{} and Memory-only (95\% interval $[-10.4,-2.2]$). The result suggests that
better reconstruction alone does not explain reliability: exposing the residual error to
the controller changes which footholds are selected.

\subsection{Calibration and Selective Planning}

Before temperature scaling, the risk head has ECE 0.091 and is overconfident on cells older
than 0.6~s. Scaling reduces ECE to 0.036 and Brier score from 0.118 to 0.082. When the planner
rejects the riskiest 20\% of candidate footholds, the remaining set contains 91\% of valid
contacts and only 24\% of invalid contacts. This selective behavior matters more than raw
height MAE: two methods can reconstruct similar surfaces while assigning very different
costs to unobserved gap edges.

\begin{table}[t]
\centering
\caption{Ablation under structured dropout. ``Contact'' applies
Eq.~\eqref{eq:contact}; ``cal.'' applies temperature calibration.}
\label{tab:ablation}
\small
\setlength{\tabcolsep}{4.2pt}
\begin{tabular}{lccc}
\toprule
Variant & Success $\uparrow$ & Interv. $\downarrow$ & Speed $\uparrow$ \\
\midrule
Full \method{}              & \best{90.8} & \best{7.9} & 0.92 \\
$-$ contact correction      & 86.1 & 10.7 & 0.91 \\
$-$ calibrated risk         & 82.9 & 13.2 & \best{0.95} \\
$-$ temporal memory         & 76.8 & 16.9 & 0.93 \\
$-$ fallback behavior       & 84.0 & 12.6 & 0.94 \\
\bottomrule
\end{tabular}
\end{table}

\subsection{Ablations}

Table~\ref{tab:ablation} separates representation and control choices. Removing temporal
memory causes the largest success loss because the planner repeatedly encounters empty
foothold sets. Removing calibration causes the highest speed but nearly doubles
interventions: unscaled scores are too confident under long occlusions. Contact correction
primarily helps after a foot lands on gravel or foam, when visual height is biased. The
fallback affects only 8.3\% of replans yet prevents five falls, indicating that explicit
behavior for ``nothing is trustworthy'' is valuable.
